% ============================================================================
% preamble.tex — 《微积分物理学》包、字体、排版、数学/物理环境
% ============================================================================

% --- 中文支持 (XeLaTeX) ---
\usepackage[fontset=fandol]{ctex}  % 中文排版核心，使用 Fandol 字体集
\usepackage{fontspec}              % 字体选择

% --- 页面布局 ---
\usepackage[
  a4paper,
  top=2.5cm, bottom=2.5cm,
  left=2.5cm, right=2.5cm
]{geometry}

% --- 数学 ---
\usepackage{amsmath, amssymb, amsthm}
\usepackage{mathrsfs}          % \mathscr 花体
\usepackage{mathtools}         % dcases, coloneqq 等
\usepackage{bm}                % 粗体数学符号
\usepackage{esint}             % 闭合曲面积分符号 \oiint
\usepackage{physics}           % 物理符号: \vb, \dd, \dv, \pdv, \grad, \div, \curl

% --- 图表 ---
\usepackage{graphicx}
\usepackage{float}
\usepackage{booktabs}
\usepackage{longtable}
\usepackage{tabularx}

% --- 绘图 ---
\usepackage{tikz}
\usetikzlibrary{
  positioning, calc, arrows.meta,
  decorations.pathreplacing,
  decorations.pathmorphing,
  decorations.markings,
  patterns,
  angles, quotes,
}
\usepackage[american]{circuitikz}  % 电路图
\usepackage{pgfplots}
\pgfplotsset{compat=1.18}
\usepgfplotslibrary{fillbetween}

% --- 交叉引用与超链接 ---
\usepackage{hyperref}
\hypersetup{
  colorlinks=true,
  linkcolor=blue!60!black,
  citecolor=green!40!black,
  urlcolor=cyan!50!black,
}
\usepackage[nameinlink]{cleveref}  % 智能引用: \cref

% --- 索引 ---
\usepackage{makeidx}
\makeindex

% --- 颜色 ---
\usepackage{xcolor}

% --- 彩色定理框 ---
\usepackage[most]{tcolorbox}
\tcbuselibrary{breakable, skins}

% ============================================================================
% 定理环境
% ============================================================================

\theoremstyle{plain}
\newtheorem{theorem}{定理}[chapter]
\newtheorem{proposition}[theorem]{命题}
\newtheorem{lemma}[theorem]{引理}
\newtheorem{corollary}[theorem]{推论}
\newtheorem{law}[theorem]{定律}

\theoremstyle{definition}
\newtheorem{definition}[theorem]{定义}
\newtheorem{example}[theorem]{例题}
\newtheorem{exercise}{习题}[chapter]

\theoremstyle{remark}
\newtheorem{remark}[theorem]{注记}

% --- tcolorbox 彩色框 ---
% 定理/命题/引理/推论/定律：蓝色
\tcolorboxenvironment{theorem}{
  enhanced, breakable, arc=3pt,
  before skip=10pt, after skip=10pt,
  colback=blue!4, colframe=blue!40!black,
  boxrule=0.6pt, left=6pt, right=6pt, top=6pt, bottom=6pt,
}
\tcolorboxenvironment{proposition}{
  enhanced, breakable, arc=3pt,
  before skip=10pt, after skip=10pt,
  colback=blue!4, colframe=blue!35!black,
  boxrule=0.6pt, left=6pt, right=6pt, top=6pt, bottom=6pt,
}
\tcolorboxenvironment{lemma}{
  enhanced, breakable, arc=3pt,
  before skip=10pt, after skip=10pt,
  colback=blue!4, colframe=blue!35!black,
  boxrule=0.6pt, left=6pt, right=6pt, top=6pt, bottom=6pt,
}
\tcolorboxenvironment{corollary}{
  enhanced, breakable, arc=3pt,
  before skip=10pt, after skip=10pt,
  colback=blue!4, colframe=blue!35!black,
  boxrule=0.6pt, left=6pt, right=6pt, top=6pt, bottom=6pt,
}
\tcolorboxenvironment{law}{
  enhanced, breakable, arc=3pt,
  before skip=10pt, after skip=10pt,
  colback=blue!6, colframe=blue!50!black,
  boxrule=0.8pt, left=6pt, right=6pt, top=6pt, bottom=6pt,
}

% 定义：绿色
\tcolorboxenvironment{definition}{
  enhanced, breakable, arc=3pt,
  before skip=10pt, after skip=10pt,
  colback=green!4, colframe=green!40!black,
  boxrule=0.6pt, left=6pt, right=6pt, top=6pt, bottom=6pt,
}

% 例题：橙色
\tcolorboxenvironment{example}{
  enhanced, breakable, arc=3pt,
  before skip=10pt, after skip=10pt,
  colback=orange!4, colframe=orange!50!black,
  boxrule=0.6pt, left=6pt, right=6pt, top=6pt, bottom=6pt,
}

% 注记：灰色
\tcolorboxenvironment{remark}{
  enhanced, breakable, arc=3pt,
  before skip=10pt, after skip=10pt,
  colback=black!3, colframe=black!30,
  boxrule=0.5pt, left=6pt, right=6pt, top=6pt, bottom=6pt,
}


% ============================================================================
% 物理常用符号
% ============================================================================
\newcommand{\vect}[1]{\vb{#1}}          % 矢量
\newcommand{\uvec}[1]{\vu{#1}}          % 单位矢量
\newcommand{\magn}[1]{\left|#1\right|}  % 模
\newcommand{\avg}[1]{\langle #1 \rangle} % 平均值
% \DeclareSIUnit\lightyear{ly}  % 需要 siunitx 包，暂不使用
